\subsection{Ausführbare Ergebnisse der Sitzung}
Für die Fortsetzung vom 29.~September tragen außerdem diese Originalstellen die native Anschlussfrage:
\begin{itemize}
\item \file{experiments/theory-contracts/raw-carrier-origin-gate-20260920/PROOF.txt}, ab Zeile 136: ursprünglicher reeller Träger $E_{\R}\oplus V_6$.
\item \file{experiments/theory-contracts/source-joint-spin-family-lift-20260922/PROOF.md}, ab Zeile 38: gemeinsame Liftfamilie und ihre Zentrumsglieder.
\item \file{experiments/theory-contracts/compiler-source-channel-gate-20260918/PROOF.txt}, ab Zeilen 133 und 167: Unterschied zwischen geometrischem Deckoperator und Glue-Clock sowie Erweiterungsindizes.
\item \file{_archive/tfpt-45/source_extracts/01_boundary_kernel_source.tex}, Zeilen 65--170 sowie 493 und 637: eingesetzte Dirac-Collarform, geometrische/Calderón-Reflexion und Übergang von elliptischem zu Dirac-Typ.
\item \file{verification/v155_quasifree_boundary.py}, ab Zeile 126; \file{v156_seam_net_construction.py}, ab Zeile 63; \file{v157_rigid_fixed_point.py}, ab Zeile 66; \file{v160_seam_gaussianity_from_pf.py}, ab Zeile 257; \file{v161_one_particle_reduction.py}, ab Zeile 170: freier Eingang, skalarer DtN-Wert, volle Schwinger-Primitivität und noch offene Operatoridentifikation. Die verkürzten Dateinamen liegen alle unter \file{verification/}.
\item \file{verification/v486_transfer_full_rule.py} und \file{verification/v814_k5_sixstep_transport.py}: ursprüngliche Clockregel, ihr Dreieroperator und die Reichweite der Sechsschrittidentifikation.
\end{itemize}
Die neuen ausführbaren Formeln stehen in \file{tfpt_explorer/source_selection.py::native_spin_anomaly} und \file{tfpt_explorer/source_realization.py::two_interval_replica_response}. Die bestehenden Builder geben sie gemeinsam mit den bisherigen Quellrechnungen aus. Die ältere Browser-Tour übernimmt diese Fortsetzung nicht automatisch.

Der neue Anhang ist zusätzlich in \file{source_realization.py::replica_pairing_readout},
\file{marked_source_process.py::replica_luders_action} und
\file{marked_source_process.py::replica_recorded_events} umgesetzt.
Die Funktionen verwenden die bestehenden Wurzeln und markierten Reflexionen. Die P1-Präzisierung wurde außerdem gegen \file{verification/v258_dirac_covariance_induction.py} und die Originalverträge \file{source-variation-origin-gate-20260922} sowie \file{source-boundary-lift-selection-20260922} geprüft. Dort stehen eine bedingte Kovarianzrekonstruktion beziehungsweise die offen bleibende Auswahl von Verbindung und Operator, kein schon bewiesener primitiver Dirac-Selektor.

Die jüngste Fortsetzung ergänzt \file{source_realization.py::anchor_clock_readout},
\file{anchor_clock_channel} und \file{family_charge_fibre}, außerdem
\file{marked_source_process.py::native_trine_dilation} und
\file{source_selection.py::affine_sewing_exclusion}.
Der markierte Intertwiner benutzt die Charakterformeln aus Abschnitt 4 des Originalvertrags \file{compiler-source-channel-gate-20260918/PROOF.txt}.
\file{verification/v221_seam_qecc.py} setzt die historischen sechsten Potenzen und die Kontrastrichtungen ein; \file{v486_transfer_full_rule.py} hält die Richtung Spektrum zu Raten ausdrücklich offen. Ihre Ausgaben wurden nicht als bereits bewiesene Auswahl der neuen Detektorbasis umgedeutet.

Die erneute Gesamtprüfung arbeitet analytisch an den tragenden Verbindungen; sie fügt keinen weiteren endlichen Quellen-Demozweig hinzu. Dafür wurden insbesondere die folgenden Originaleingänge gegeneinander gelesen:
\begin{itemize}
\item \file{_archive/tfpt-45/source_extracts/01_boundary_kernel_source.tex}, um Zeilen 470, 544, 629 und 937: Eingang der operativen Kategorie, Spektralfluss, Defektminimum und relativer APS-Rahmen.
\item \file{_archive/tfpt-45/source_extracts/02_carrier_source.tex}, um Zeile 3005: Mastervariation über bereits gewähltem diskretem Datum.
\item \file{_archive/tfpt-45/source_extracts/04_qft_source.tex}, um Zeilen 1750, 2150 und 2740: PhysAdm, effektive Wirkung und bedingte OS-/Streurekonstruktion.
\item \file{verification/v154_simple_current_theorem.py}, um Zeilen 53 und 95: exakte Erweiterungsarithmetik und externer physischer Netz-Eingang.
\item \file{verification/v260_k3_kummer_unification.py}, um Zeile 109; \file{verification/v367_seam_s3_lattice.py}, um Zeile 118; \file{verification/v470_alpha_inflow_level.py}, um Zeilen 125 und 161: K3, eingesetzte Majoranazahl und Unterschied zwischen Ladungs- und thermischem Index.
\item \file{tfpt_research_contracts.tex}, um Zeilen 1394, 12698, 13235, 13616 und 13663: ursprüngliche Grenzen von K3, Prime Completion, Masterwirkung, Determinantenlinie und gemeinsamem erzeugendem Funktional.
\item \file{experiments/theory-contracts/source-variation-origin-gate-20260922/PROOF.md} und \file{source-three-route-closure-20260922/SELECTION.md} im selben Vertragsverzeichnis: bereits dokumentierte Herkunfts- und Eindeutigkeitsgrenzen.
\end{itemize}
Die neuen analytischen Anschlüsse wurden unabhängig auf Gruppen, Grade, Statistik und Zustandsvoraussetzungen gegengeprüft. Sie sind hier mit Beweisen beziehungsweise Primärsätzen angegeben; eine neue Lean-Formalisierung oder eine vollständige erneute Ausführung aller Repo-Verifikationen wird damit nicht behauptet.

Die Herkunfts- und Transferprüfung der Fortsetzung verwendet zusätzlich folgende konkrete Originalstellen:
\begin{itemize}
\item \file{experiments/theory-contracts/microscopic-charged-car-limit/README.md}, Zeilen 167--202: ursprüngliche Viertelholonomie-Polarisation und geladene CAR-Moden.
\item \file{experiments/theory-contracts/microscopic-neutral-limit/README.md}, Abschnitte 5--7, und \file{microscopic-fourpoint-limit/README.md}, Abschnitte 3--4 im selben Vertragsverzeichnis: vollständige bewegte Quellgeschichten und logarithmischer Hilfsfenstervergleich. Abschnitt~\ref{sec:microscopicpairbridge} führt daraus den neuen positiven distributionswertigen Achtkopien-Paarkern her. Er ist schriftlich bewiesen und unabhängig mathematisch gegengeprüft, nicht in Lean formalisiert.
\item \file{experiments/theory-contracts/compiler-source-channel-gate-20260918/PROOF.txt}, Zeilen 27--83: tatsächliche Kanalzahl und Rang eins der führenden Strom-Grammatrix.
\item \file{verification/v498_celestial_wp5b_singular_vector.py}, Abschnitte S4.10 und S4.13: Verhältnis von Viertel-Slot- und Periodenmoden sowie der bereits gitterbedingte Nullvektor.
\item \file{verification/v258_dirac_covariance_induction.py}, Präambel und Schlussprüfung: $\mu_{\rm geo}\Pi_{\rm odd}H_{\rm mod}$ und bedingter physischer Kompressionspfeil.
\item \file{experiments/theory-contracts/source-yukawa-forward-gate-20260921/yukawa_source_provenance_20260921.md}, ab Zeile 143: bereits geprüfte Vorwärtsschranken der Kovarianzroute.
\item \file{tfpt_research_contracts.tex}, Kennungen \file{TFPT4D.LATTICE.ACTION.01}, \file{SEAM.DETLINE.UNIFICATION.01} und \file{FTRANSFER.GENERATING.01}: gemeinsamer physischer Abschlussvertrag.
\end{itemize}

Alle folgenden Dateien liegen im Repository unter \file{tfpt_explorer/}. Sie sind der unmittelbare Ursprung der hier beschriebenen neuen endlichen Rechnungen; ihre Existenz ist kein zusätzlicher Beweis über ihren ausgewiesenen Geltungsbereich hinaus.
\begin{longtable}{L{0.37\linewidth}L{0.55\linewidth}}
\toprule
Datei & Inhalt\\\midrule\endhead
\file{origin_closure.py}, \file{core.py} & Originalanker, Rückkopplungen, ältere TFPT-Auslesungen\\
\file{process.py} & Ereignisse, Quartikcode, Bindung, Dreiecksrekursion\\
\file{quartic_selection.py}, \file{consolidation.py} & Molien-Auswahl, Präparationsfläche, Projektoren und Kontrollen\\
\file{composition.py} & Pfadblöcke, Neun-Port-Komposition und Leakage\\
\file{joint_constraints.py}, \file{spatial_response.py} & Gemeinsame Bindungsform und volle räumliche Antwort\\
\file{origin_transfer.py} & Ursprünglicher Transfer, Klasse und Fortsetzung\\
\file{cartan_source.py}, \file{marked_source_process.py} & Markierungsprüfung, Fixalgebren, Trine und native Ereignisfaser\\
\file{history_kernel.py} & Positiver kohärenter Geschichtskern und Rangvergleich\\
\file{hecke_source.py} & Untergitter, Charakterfortsetzung, Sektorprojektoren, geladener Carry\\
\file{twisted_source.py}, \file{torus_lift.py}, \file{source_process.py} & Tatsächliche Lifts, Phasen und ausgeschlossene Gleichsetzungen\\
\file{current_block_geometry.py} & $C_4+X_{20}$, Chevalley-Klammern, gemeinsamer Quellabschluss\\
\file{rewrite_coherence.py} & KZ-/Ward-Kompositionskonsistenz\\
\file{source_unification.py} & $A_8$-Gitter, $\Z_3$-Erweiterung und gemeinsamer Virasorovektor\\
\file{source_program.py} & Ganze Stromwörter, verbundene Antworten, harmonischer $W$-Anschluss\\
\file{source_selection.py} & Allgemeiner affiner Vierpunkttensor, Levelselektoren und explizite Level-zwei-Kontrolle\\
\file{source_realization.py} & Originale SU9-Feldabbildung, unabhängiger X-Wickweg, voller C/X-Gitterauswerter, positiver Geschichtskernel, rationale Momente und Möbius-Nachfahrenturm\\
\file{README.md} & Bedienung, Befunde und ausdrückliche Grenzen des Explorers\\\bottomrule
\end{longtable}

Die wichtigen Laufprotokolle liegen unter \file{tfpt_explorer/runtime/}:
\file{snapshot.json},
\file{verification_run_summary.json},
\file{verification_full.log},
\file{build_audit_summary.json},
\file{lean_replay.json},
\file{browser_check_summary.json},
\file{phase_selection_review.json}.
Der Stand der einzelnen Dateien unterscheidet sich zeitlich. Beispielsweise ist die frühe Browserzusammenfassung kein Nachweis aller später hinzugefügten Tourstationen.

Der exportierte vollständige Ablauf liegt unter
\file{/Users/stefanhamann/Documents/Codex/2026-09-28/unter/outputs/TFPT_Ablauf.json}.
Die frühere HTML-Gesamtprüfung,
\file{TFPT_Gesamtpruefung_2026-09-28.html},
und \file{Quellenfamilie_Pruefergebnis.json} im selben Ausgabeordner dokumentieren einen älteren Zwischenstand. Dieses LaTeX-Dokument konsolidiert die danach hinzugekommenen Verbindungen und Korrekturen.

\subsection{Historischer Algorithmusvergleich}
Der im Anwendungsteil verwendete Vergleich ist am datierten Original nachgelesen:
\file{/Users/stefanhamann/Documents/Codex/2026-09-25/und-verifiziere-und-beweise-und-teste/BEURTEILUNG.md}.
Er umfasst den reparierten endlichen Causal-Compiler, AALpy-Vergleiche, unabhängige Modellprüfungen, SQLite-Traces und Gegenfälle. Er belegt weder eine Hylæan-Integration noch eine Überlegenheit gegenüber modernen neuronalen Weltmodellen. Er ist hier aufgenommen, damit dieselbe bereits geprüfte Idee nicht erneut als ungetestete Neuheit erscheint.

\section{Begriffe in einfacher Sprache}
\begin{longtable}{L{0.22\linewidth}L{0.70\linewidth}}
\toprule
Begriff & Gemeinte Bedeutung\\\midrule\endhead
Quelle & Algebra von Operationen samt Zustand, aus denen gemeinsame Antworten berechnet werden\\
Ladung & Kennzeichnung, wie ein Feld unter einer bestimmten Symmetrie transformiert\\
Kokzyklus & Konsistente Phasenregel, die beim Zusammensetzen geladener Operationen benötigt wird\\
Carry & Mitgeführte Information, die bei späteren Operationen wieder beobachtbar werden kann\\
Intertwiner & Abbildung, bei der ``erst operieren, dann übersetzen'' gleich ``erst übersetzen, dann operieren'' ist\\
Ward-Regel & Rekursion, die höhere Stromantworten durch Paarung und Lieklammer bestimmt\\
Verbunden & Der Teil einer Antwort, der nach Abzug der Produkte echter Teilantworten übrig bleibt\\
Hyperkante & Gemeinsame Antwort mehrerer Einsetzungen; im Dokument zunächst ein Rechen- und Antwortobjekt\\
GNS & Aufbau eines Hilbertraums aus einer bereits gegebenen positiven Quelle\\
Hankelmatrix & Tabelle aller Antworten von Vergangenheit gefolgt von Zukunft; ihr Rang misst lineare Vorhersagedimension\\
Vergröberung & Wechsel der Auflösung; je nach Definition entweder antworttreue Rekodierung oder echtes Vergessen\\
Fixpunkt & Objekt, das unter einer festgelegten Abbildung unverändert bleibt; Gesetz und Zustand getrennt betrachten\\
Virasorovektor & Das algebraische Objekt, das den gemeinsamen konformen Stress-Tensor und seine Zeit-/Skalenwirkung bestimmt\\
KZ-Verbindung & Regel für konsistenten Transport konformer Blöcke beim Verschieben ihrer Einsetzungen\\
Nachfahren & Weitere Zustände/Felder derselben Quelle, erzeugt durch ihre Strommoden\\
Markierung & Physisch oder geometrisch ausgezeichnete Struktur, die bei einer angeblich gleichen Darstellung erhalten werden muss\\
Gesamtlösung & Gemeinsame Herkunft, Dynamik, Raumzeitabbildung und überprüfbare Antworten derselben Konstruktion\\\bottomrule
\end{longtable}

\begin{thebibliography}{99}
\bibitem{tomiyamaexpectation}
Jun Tomiyama: \emph{On the projection of norm one in $W^*$-algebras}.
Proceedings of the Japan Academy 33 (1957), 608--612.
\url{https://doi.org/10.3792/pja/1195524885}.
Die Originalaussage behandelt Positivität und Bimodulstruktur; sie ersetzt keine allgemeine Produktregel.

\bibitem{pollockprocess}
Felix A. Pollock, Cesar Rodriguez-Rosario, Thomas Frauenheim, Mauro Paternostro und Kavan Modi:
\emph{Non-Markovian quantum processes: Complete framework and efficient characterization}.
Physical Review A 97 (2018), 012127.
\url{https://arxiv.org/abs/1512.00589}.
Ein Rahmen für vollständige Mehrzeitantworten mit Gedächtnis, kein TFPT-Ursprungsbeweis.

\bibitem{dongmasonlow} C. Dong, G. Mason: \emph{Holomorphic Vertex Operator Algebras of Small Central Charges}. Klassifikation stark rationaler holomorpher VOAs mit $c=8$ und $c=16$ als entsprechende gerade selbstduale Gittertheorien. \url{https://arxiv.org/abs/math/0203005}.
\bibitem{criadoredefinitions} J. C. Criado, M. P\'erez-Victoria: \emph{Field redefinitions in effective theories at higher orders}. JHEP 03 (2019), 038. Physische Observablen versus Feldpräsentation; korrekter Transport der Quellen und höhere Ordnungen. \url{https://arxiv.org/abs/1811.09413}.
\bibitem{nikulinforms} V. V. Nikulin: \emph{Integral symmetric bilinear forms and some of their applications}. Mathematics of the USSR-Izvestiya 14 (1980), 103--167. Klassische Diskriminantenformen- und Überlagerungsgittertechnik; die globale Anwendung auf $A_{d-1}\oplus D_{d+1}$ ist in Abschnitt~\ref{sec:globaloriginselection} ausgeschrieben. \url{https://doi.org/10.1070/IM1980v014n01ABEH001060}.
\addcontentsline{toc}{section}{Wissenschaftliche Primärliteratur}
\bibitem{razamattong} S. S. Razamat, D. Tong: \emph{Gapped Chiral Fermions}. Abschnitt 3, Gleichung (7), zur Spin-$\Z_4$-Anomalie. \url{https://arxiv.org/html/2009.05037v3}.
\bibitem{garciamontero} I. García-Etxebarria, M. Montero: \emph{Dai-Freed anomalies in particle physics}. Abschnitt 4.3 zur Spin-$\Z_4$-Struktur, dem rechtshändigen Neutrino und Majoranamassen. \url{https://arxiv.org/html/1808.00009v3}.
\bibitem{boylesmith} P. Boyle Smith, Y.-H. Lin, Y. Tachikawa, Y. Zheng: \emph{Classification of chiral fermionic CFTs of central charge $\leq16$}. SciPost Physics 16 (2024), 058. Abschnitt 5.1, Gleichungen (5.3)--(5.6). \url{https://arxiv.org/abs/2303.16917}.
\bibitem{beemchiral} C. Beem, M. Lemos, P. Liendo, W. Peelaers, L. Rastelli, B. C. van Rees: \emph{Infinite Chiral Symmetry in Four Dimensions}. Abschnitt 3 und Gleichung (57). \url{https://arxiv.org/html/1312.5344v3}.
\bibitem{costellopaquette} K. Costello, N. M. Paquette: \emph{Associativity of One-Loop Corrections to the Celestial Operator Product Expansion}. Insbesondere die Strom-OPE und die für die Anomaliekompensation zugelassenen Gruppen. \url{https://arxiv.org/html/2204.05301v1}.
\bibitem{arakidomain} S. Araki, H. Fukaya, T. Onogi, S. Yamaguchi: \emph{Symmetric Mass Generation for Domain-Wall Fermions}. Zweidimensionale euklidische Gitterrechnung mit acht Majorana-Flavors. \url{https://arxiv.org/html/2608.29963v1}.
\bibitem{hasenfratzguide} A. Hasenfratz, C. Xu: \emph{A Guide to Symmetric Mass Generation in Lattice-QCD}. Fassung 3, zur Bedeutung der konkreten UV-Wechselwirkung. \url{https://arxiv.org/html/2604.02424v3}.
\bibitem{gavrylenkomarshakov} P. Gavrylenko, A. Marshakov: \emph{Free fermions, W-algebras and isomonodromic deformations}. Abschnitte 4.1--5 zu Monodromiedefekten, Fermionenkernel und Determinanten. \url{https://arxiv.org/abs/1605.04554}.
\bibitem{bfkoriginal} D. Burghelea, L. Friedlander, T. Kappeler: \emph{Meyer--Vietoris type formula for determinants of elliptic differential operators}. Journal of Functional Analysis 107 (1992), 34--65. \url{https://doi.org/10.1016/0022-1236(92)90099-5}.
\bibitem{leebfk} Y. Lee: \emph{Burghelea--Friedlander--Kappeler's gluing formula for the zeta-determinant and its applications}. Zu lokalen Verklebungsfaktoren unter Produktstruktur am Rand. \url{https://arxiv.org/abs/math/0304250}.
\bibitem{yangmodular} T. H. Yang: \emph{Modular commutator as a robust topological invariant and approximate Markovianity}. 8.~September 2026, Theoreme 2.9 und 3.11 sowie Abschnitt 5.2.1. \url{https://arxiv.org/abs/2609.09019}.
\bibitem{gasslevinreplica} J. Gass, M. Levin: \emph{R\'enyi-like entanglement probe of the chiral central charge}. Physical Review B, 22.~Juni 2026; Gleichungen (5), (10)--(14) und ausdrücklich benannte Geltungsgrenzen. \url{https://arxiv.org/abs/2512.20608}.
\bibitem{shefferdefects} Y. Sheffer, R. Fan, A. Stern, E. Berg, S. Ryu: \emph{Probing chiral topological states with permutation defects}. Physical Review B, 15.~Mai 2026; Abschnitte II--V und Anhang I. \url{https://arxiv.org/abs/2512.04649}.
\bibitem{castroentropy} A. Castro, S. Detournay, N. Iqbal, E. Perlmutter: \emph{Holographic entanglement entropy and gravitational anomalies}. Gleichungen (25), (28) zur Rand-CFT mit festgehaltenem Regulatorrahmen. \url{https://arxiv.org/abs/1405.2792}.
\bibitem{thorngrendisentanglers} R. Thorngren, J. Preskill, L. Fidkowski: \emph{Chiral lattice gauge theories from symmetry disentanglers}. JHEP 07 (2026), 271, Abschnitte 1.1 und 5. \url{https://arxiv.org/abs/2601.04304}.
\bibitem{kapustinbordism} A. Kapustin, R. Thorngren, A. Turzillo, Z. Wang: \emph{Fermionic Symmetry Protected Topological Phases and Cobordisms}. Anhang, Gleichung (6), zur Quantisierung der gravitativen Chern--Simons-Antwort auf Spin- und orientierten Hintergründen. \url{https://arxiv.org/abs/1406.7329}.
\bibitem{freedhopkins} D. S. Freed, M. J. Hopkins: \emph{Reflection positivity and invertible topological phases}. Insbesondere Abschnitt 5.4, Conjecture 8.37, Remark 8.39 und Theorem 9.79 mit Folgerungen zur Unterscheidung topologischer Klassifikation, physischer Interpretation und Symmetrietypen. \url{https://arxiv.org/abs/1604.06527}.
\bibitem{absindex} M. F. Atiyah, R. Bott, A. Shapiro: \emph{Clifford Modules}. Topology 3, Supplement 1 (1964), 3--38. Abschnitte 6 und 11 zu reellen/komplexen Indexnormalisierungen. \url{https://www.math.ens.psl.eu/~benoist/refs/ABS.pdf}.
\bibitem{atiyahsinger3} M. F. Atiyah, I. M. Singer: \emph{The Index of Elliptic Operators: III}. Annals of Mathematics 87 (1968), 546--604. \url{https://annals.math.princeton.edu/1968/87-3/p07}.
\bibitem{chuzheng} Y.-J. Chu, Z.-J. Zheng: \emph{Vertex Operator Algebra Analogue of Embedding $D_8$ into $E_8$}. \url{https://arxiv.org/abs/0808.1458}.
\bibitem{barronfermion} K. Barron, N. Vander Werf: \emph{On permutation-twisted free fermions and two conjectures}. Abschnitt 6, Gleichungen (6.1)--(6.2) und anschließende Tensorpotenz-Isomorphie $V_{\Z^d}\cong\mathcal F^{\otimes2d}$. \url{https://arxiv.org/abs/1310.1958}.
\bibitem{bockenhauercar} J. B\"ockenhau-er: \emph{An Algebraic Formulation of Level One Wess--Zumino--Witten Models}. Abschnitte 2--4 zu den $\mathfrak{so}(N)_1$-Sektoren, NS-Moden und Virasorogeneratoren. \url{https://arxiv.org/abs/hep-th/9507047}.
\bibitem{dongxu} C. Dong, F. Xu: \emph{Conformal Nets Associated with Lattices and Their Orbifolds}. Proposition 3.15, Corollary 3.19 und Lemma 4.1 zu Gitternetzsektoren und internen Automorphismen. \url{https://arxiv.org/abs/math/0411499}.
\bibitem{doplicherroberts} S. Doplicher, J. E. Roberts: \emph{Endomorphisms of $C^*$-algebras, cross products and duality for compact groups}. Annals of Mathematics 130 (1989), 75--119. \url{https://annals.math.princeton.edu/1989/130-1/p03}.
\bibitem{rehren} K.-H. Rehren: \emph{Algebraic Holography}. Corollaries 1--2 und Proposition zum chiralen Rand und AdS$_2$. \url{https://arxiv.org/abs/hep-th/9905179}.
\bibitem{bdfs} D. Buchholz, O. Dreyer, M. Florig, S. J. Summers: \emph{Geometric Modular Action and Spacetime Symmetry Groups}. Theoreme 4.3.9 und 5.1.2 sowie Proposition 5.1.3; die Keilordnung gehört zum Eingang. \url{https://arxiv.org/abs/math-ph/9805026}.
\bibitem{schroer} B. Schroer: \emph{Lightfront Formalism versus Holography and Chiral Scanning}. Abschnitt 4 und Anhang zur zusätzlichen transversalen Lokalisation. \url{https://arxiv.org/abs/hep-th/0108203}.
\bibitem{os1975} K. Osterwalder, R. Schrader: \emph{Axioms for Euclidean Green's Functions II}. Communications in Mathematical Physics 42 (1975), 281--305. \url{https://doi.org/10.1007/BF01608978}.
\bibitem{gorardcausal} J. Gorard: \emph{Some Relativistic and Gravitational Properties of the Wolfram Model}. \url{https://arxiv.org/abs/2004.14810}; sowie \emph{Algorithmic Causal Sets and the Wolfram Model}. \url{https://arxiv.org/abs/2011.12174}. Verwendet zur Trennung kausaler Konsistenz von zusätzlichen Dimensions- und Dynamikannahmen.
\bibitem{cklw} S. Carpi, Y. Kawahigashi, R. Longo, M. Weiner: \emph{From vertex operator algebras to conformal nets and back}. \url{https://arxiv.org/abs/1503.01260}. Verwendet für den bedingten Übergang unitärer stark lokaler VOAs zu konformen Netzen.
\bibitem{donglin} C. Dong, X. Lin: \emph{Unitary vertex operator algebras}. Insbesondere Sätze 4.7--4.8 zur unitären affinen Unteralgebra. \url{https://arxiv.org/abs/1308.2361}.
\bibitem{frenkelkac} I. B. Frenkel, V. G. Kac: \emph{Basic representations of affine Lie algebras and dual resonance models}. Inventiones mathematicae 62 (1980), 23--66. \url{https://doi.org/10.1007/BF01391662}. Grundlegende Gittervertex-Realisierung; die Formeln der Fortsetzung werden mit den ursprünglichen TFPT-Wurzeln und Kokzyklen konkret ausgewertet.
\bibitem{laszlosorger} Y. Laszlo, C. Sorger: \emph{The line bundles on the stack of parabolic G-bundles over curves and their sections}. Abschnitt 2.3--2.6 zum Dynkinindex und Determinantenkokzyklus. \url{https://arxiv.org/abs/alg-geom/9507002}.
\bibitem{perelomovstates} D. Ojeda-Guillén, R. D. Mota, V. D. Granados: \emph{SU(1,1) and SU(2) Perelomov number coherent states: algebraic approach for general systems}. \url{https://arxiv.org/abs/1206.1555}. Die hier verwendete $h=1$-Reihe wird zusätzlich unmittelbar aus den Strommoden hergeleitet.
\bibitem{longomodular} R. Longo, P. Martinetti, K.-H. Rehren: \emph{Geometric modular action for disjoint intervals and boundary conformal field theory}. \url{https://arxiv.org/abs/0912.1106}. Mehrintervall-Modularstruktur und ihre Abgrenzung von der einfachen Möbiusbewegung eines Intervalls.
\bibitem{distler} J. Distler, E. Sharpe: \emph{Heterotic compactifications with principal bundles for general groups and general levels}. Insbesondere Abschnitt 6.1 zur $SU(9)/\Z_3$-Erweiterung. \url{https://arxiv.org/abs/hep-th/0701244}.
\bibitem{gui} B. Gui: \emph{Convergence of Sewing Conformal Blocks}. \url{https://arxiv.org/abs/2011.07450}. Analytische Voraussetzungen für Sewing.
\bibitem{guizhang} B. Gui, H. Zhang: \emph{Analytic Conformal Blocks of $C_2$-cofinite Vertex Operator Algebras III: The Sewing-Factorization Theorems}. Proc. Lond. Math. Soc. 132 (2026), e70130. \url{https://arxiv.org/abs/2503.23995}. Kein Ersatz für die Auswahl der physischen TFPT-Quelle.
\bibitem{carlyle} J. W. Carlyle, A. Paz: \emph{Realizations by stochastic finite automata}. Journal of Computer and System Sciences 5 (1971), 26--40. \url{https://www.sciencedirect.com/science/article/pii/S0022000071800053}.
\bibitem{segal} I. E. Segal: \emph{Irreducible representations of operator algebras}. Bulletin of the AMS 53 (1947). \url{https://www.ams.org/bull/1947-53-02/S0002-9904-1947-08742-5/S0002-9904-1947-08742-5.pdf}.
\bibitem{combs} G. Chiribella, G. M. D'Ariano, P. Perinotti: \emph{Theoretical framework for quantum networks}. Physical Review A 80, 022339 (2009). \url{https://arxiv.org/abs/0904.4483}.
\bibitem{zuse} K. Zuse: \emph{Rechnender Raum} (1969). Originalscan: \url{https://sferics.idsia.ch/pub/juergen/zuserechnenderraum.pdf}.
\bibitem{wolfram} S. Wolfram: \emph{A Class of Models with the Potential to Represent Fundamental Physics}, Abschnitt 5.19, \emph{Weighted Multiway Graphs}. \url{https://www.wolframphysics.org/technical-introduction/the-updating-process-for-string-substitution-systems/weighted-multiway-graphs/}.
\bibitem{littman} M. L. Littman, R. S. Sutton, S. Singh: \emph{Predictive Representations of State} (2001). \url{https://proceedings.neurips.cc/paper/2001/file/1e4d36177d71bbb3558e43af9577d70e-Paper.pdf}.
\bibitem{thon} M. Thon, H. Jaeger: \emph{Links between multiplicity automata, observable operator models and predictive state representations: a unified learning framework}. JMLR 16 (2015). \url{https://www.jmlr.org/papers/v16/thon15a.html}.
\bibitem{adhikary} S. Adhikary et al.: \emph{Quantum Tensor Networks, Stochastic Processes, and Weighted Automata}. AISTATS 2021. \url{https://proceedings.mlr.press/v130/adhikary21a.html}. Verwendet als Neuheitsvergleich der Modellklassen.
\bibitem{lsharp} F. Vaandrager, B. Garhewal, J. Rot, T. Wißmann: \emph{A New Approach for Active Automata Learning Based on Apartness}. \url{https://arxiv.org/abs/2107.05419}.
\bibitem{egg} M. Willsey et al.: \emph{egg: Fast and Extensible Equality Saturation}. \url{https://arxiv.org/abs/2004.03082}.
\bibitem{singh} S. Singh, R. N. C. Pfeifer, G. Vidal: \emph{Tensor network decompositions in the presence of a global symmetry}. \url{https://arxiv.org/abs/0907.2994}.
\bibitem{singhvidal} S. Singh, G. Vidal: \emph{Tensor network states and algorithms in the presence of a global SU(2) symmetry}. \url{https://arxiv.org/abs/1208.3919}.
\bibitem{domino} K. Domino, P. Gawron, Ł. Pawela: \emph{Efficient computation of higher order cumulant tensors}. \url{https://arxiv.org/abs/1701.05420}.
\bibitem{schneidman2003} E. Schneidman et al.: \emph{Network information and connected correlations}. \url{https://arxiv.org/abs/physics/0307072}.
\bibitem{schneidman2006} E. Schneidman et al.: \emph{Weak pairwise correlations imply strongly correlated network states in a neural population}. Nature 440 (2006). \url{https://www.nature.com/articles/nature04701}.
\bibitem{lean} L. de Moura, S. Ullrich: \emph{The Lean 4 Theorem Prover and Programming Language}. \url{https://lean-lang.org/papers/lean4.pdf}.
\bibitem{metacoq} M. Sozeau: \emph{Touring the MetaCoq Project}. EPTCS 337 (2021), 13--29. \url{https://arxiv.org/abs/2107.07670}.
\bibitem{freeddeterminant} D. S. Freed: \emph{Determinant Line Bundles Revisited}. Theorem 24 zur Krümmung der Determinantenlinie einer Diracfamilie. \url{https://arxiv.org/abs/dg-ga/9505002}.
\bibitem{baerballmann} C. Bär, W. Ballmann: \emph{Guide to Boundary Value Problems for Dirac-Type Operators}. Abschnitt 2 zur zusätzlichen Cliffordstruktur des Dirac-Symbols. \url{https://arxiv.org/abs/1307.3021}.
\bibitem{arakicar} H. Araki: \emph{On Quasifree States of CAR and Bogoliubov Automorphisms}. Publ. RIMS Kyoto Univ. 6 (1970/71), 385--442; Lemma 4.3, S. 390, zur Eindeutigkeit des Zustands bei gegebener Basisprojektion. \url{https://ems.press/content/serial-article-files/41673}.
\bibitem{calderonscattering} C. Guillarmou, S. Moroianu, J. Park: \emph{Bergman and Calderón projectors for Dirac operators}. Theorem 2 und Proposition 12 zur Calderón-Projektion und der zugehörigen konform kompakten Streukonstruktion. \url{https://www.imar.ro/~sergium/fisiere/calderon.pdf}.
\bibitem{takesakiexpectation} M. Takesaki: \emph{Conditional expectations in von Neumann algebras}. Journal of Functional Analysis 9 (1972), 306--321. \url{https://doi.org/10.1016/0022-1236(72)90004-3}.
\bibitem{connesradon} A. Connes: \emph{Noncommutative Geometry} (1994), Kapitel V, Abschnitt 5$\alpha$, Theorem 1(a), gedruckte S. 481: Radon--Nikodym-Kokzyklus und Zusammenhang modularer Gruppen. \url{https://alainconnes.org/wp-content/uploads/book94bigpdf.pdf}.
\bibitem{wiesbrockintersection} H.-W. Wiesbrock: \emph{Symmetries and Modular Intersections of Von Neumann Algebras}. Letters in Mathematical Physics 39 (1997), 203--212. Modulare Inklusionen und zusätzliche Reflexionsbedingungen; kein voraussetzungsloser $4D$-Rekonstruktionssatz. \url{https://doi.org/10.1023/A:1007361114049}.
\bibitem{kahlerwiesbrock} R. Kähler, H.-W. Wiesbrock: \emph{Modular theory and the reconstruction of four-dimensional quantum field theories}. Journal of Mathematical Physics 42 (2001). Die relative modulare Stellung der Algebren ist Rekonstruktionsvoraussetzung. \url{https://doi.org/10.1063/1.1327597}.
\bibitem{morinellirehren} V. Morinelli, K.-H. Rehren: \emph{Spacelike deformations: higher-helicity fields from scalar fields}. Letters in Mathematical Physics 110 (2020), 2019--2038. Abschnitte 2.1--2.4 zu Möbius-Zeit, räumlichen Generatoren und den benötigten Einteilchendarstellungen. \url{https://doi.org/10.1007/s11005-020-01294-w}.
\bibitem{koesterconformal} S. Köster: \emph{Conformal Transformations as Observables}. Letters in Mathematical Physics 61 (2002), 187--198. Abschnitt 2 und Proposition 1 zu konformer Zeit, Translation und positiver Energie. \url{https://arxiv.org/abs/math-ph/0201016}.
\bibitem{headrickreplica} M. Headrick: \emph{Entanglement Rényi entropies in holographic theories}. Gleichungen (A.1) und (4.24) zur Überlagerungsnormierung und Zweintervall-Replica-Antwort. \url{https://arxiv.org/abs/1006.0047}.
\bibitem{castroanomaly} A. Castro, S. Detournay, N. Iqbal, E. Perlmutter: \emph{Holographic entanglement entropy and gravitational anomalies}. Abschnitt II.B, insbesondere (2.16)--(2.17), zu getrennten chiralen Twistgewichten und Frameabhängigkeit. \url{https://arxiv.org/abs/1405.2792}.
\bibitem{petzrenyi} J. Kudler-Flam: \emph{Rényi Mutual Information in Quantum Field Theory}. Zur Unterscheidung der üblichen Rényi-Entropiekombination von einer relativen Rényi-Größe. \url{https://arxiv.org/abs/2211.01392}.
\bibitem{e8replica} L. A. León Andonayre, R. Poddar: \emph{Rényi entropy of single-character CFTs on the torus}. Physical Review D 112 (2025), 025011. Insbesondere Abschnitte 3.1--3.2 zur E8-Permutationsorbifold-Fusion und Charakterreduktion. Die dortige Torus-Einintervallgeometrie wird nicht mit unserer Sphären-Zweiintervallgeometrie gleichgesetzt. \url{https://arxiv.org/abs/2412.00192}.
\bibitem{collazuolmelnikov} V. Collazuol, I. V. Melnikov: \emph{A twist at infinite distance in the CHL string}. Abschnitt 4.1, Gleichungen (4.15)--(4.18), zum tatsächlichen $E_{8,2}$-Singulett-Twist mit Ising-Gewicht $1/2$. \url{https://arxiv.org/abs/2402.01606}.
\bibitem{ardonneising} E. Ardonne, G. Sierra:
\emph{Chiral correlators of the Ising conformal field theory}.
Journal of Physics A 43 (2010), 505402, Gleichung (19):
Cauchy-Pfaffian/Hafnian-Identität für beliebige gerade Punktzahlen.
\url{https://arxiv.org/abs/1008.2863}.

\bibitem{wolfjoint} M. M. Wolf, J. Eisert, T. S. Cubitt, J. I. Cirac: \emph{Assessing non-Markovian dynamics}. Physical Review Letters 101 (2008), 150402. Gleichung (2) und Anhang zum bedingten Choi-Kriterium. \url{https://arxiv.org/abs/0711.3172}.
\bibitem{gutoskijoint} G. Gutoski: \emph{Properties of Local Quantum Operations with Shared Entanglement}. Quantum Information and Computation 9 (2009), 739--764. \url{https://arxiv.org/abs/0805.2209}.
\bibitem{dongnagjoint} C. Dong, K. Nagatomo: \emph{Automorphism Groups and Twisted Modules for Lattice Vertex Operator Algebras}. Abschnitte 2.1--2.4 zu Kokzyklus, Isometrielift und Vakuumautomorphismen. \url{https://arxiv.org/abs/math/9808088}.
\bibitem{darianojo} G. M. D'Ariano, P. Perinotti: \emph{Derivation of the Dirac Equation from Principles of Information Processing}. Physical Review A 90 (2014), 062106. \url{https://arxiv.org/abs/1306.1934}.
\bibitem{tanimotojoint} Y. Tanimoto: \emph{Ground state representations of some non-rational conformal nets}. Symmetry 10 (2018), 415. Abschnitt 3.1 zum lokalen Stromnetz. \url{https://arxiv.org/abs/1807.11723}.
\bibitem{chamseddineconnes} A. H. Chamseddine, A. Connes: \emph{Noncommutative Geometry as a Framework for Unification of all Fundamental Interactions including Gravity. Part I}. Gleichungen (5.24), (5.47), (5.49) zum vollständigen gravitativen Spektralaktionskoeffizienten. \url{https://arxiv.org/abs/1004.0464}.
\bibitem{casasibarra} J. A. Casas, A. Ibarra: \emph{Oscillating neutrinos and $\mu\to e,\gamma$}. Gleichung (13) zur allgemeinen Seesaw-Faktorisierung. \url{https://arxiv.org/abs/hep-ph/0103065}.
\bibitem{davidsonibarra} S. Davidson, A. Ibarra: \emph{A lower bound on the right-handed neutrino mass from leptogenesis}. Schranke der Zerfallsasymmetrie, keine automatische Gleichsetzung mit einer PMNS-Phase. \url{https://arxiv.org/abs/hep-ph/0202239}.
\bibitem{covirouletvissani} L. Covi, E. Roulet, F. Vissani: \emph{CP violating decays in leptogenesis scenarios}. Gleichungen (4)--(8), mit Übersetzung zur Heavy-row-Konvention. \url{https://arxiv.org/abs/hep-ph/9605319}.
\bibitem{horavawittenorigin} P. Hořava, E. Witten: \emph{Heterotic and Type I String Dynamics from Eleven Dimensions}. Abschnitt 2 zur lokalen Wandanomalie. \url{https://arxiv.org/abs/hep-th/9510209}.
\bibitem{horavawittenwall} P. Hořava, E. Witten: \emph{Eleven-Dimensional Supergravity on a Manifold with Boundary}. Abschnitte 2--3 zu Wandmultiplet und Anomaliefaktorisierung. \url{https://arxiv.org/abs/hep-th/9603142}.
\bibitem{braunheterotic} V. Braun, Y.-H. He, B. A. Ovrut, T. Pantev: \emph{A Heterotic Standard Model}. Gleichungen (1)--(10): Eichbündelindex und Calabi--Yau-Topologie. \url{https://arxiv.org/abs/hep-th/0501070}.
\bibitem{braundeformation} V. Braun, P. Candelas, R. Davies, R. Donagi: \emph{The MSSM Spectrum from (0,2)-Deformations of the Heterotic Standard Embedding}. Einleitung, Abschnitt 3.1 und Anhang B. \url{https://arxiv.org/abs/1112.1097}.
\bibitem{heeckquadruple} J. Heeck, W. Rodejohann: \emph{Neutrinoless Quadruple Beta Decay}. Abschnitt II, Gleichungen (2)--(5): Ladung-zwei-Feld und zustandsabhängiger Dirac-Schutz. \url{https://arxiv.org/abs/1306.0580}.
\bibitem{liyautorsion} J. Li, S.-T. Yau: \emph{The Existence of Supersymmetric String Theory with Torsion}. Sätze 2.1 und 4.3, mit gegebener holomorpher Deformationsfamilie. \url{https://arxiv.org/abs/hep-th/0411136}.
\bibitem{dongradical} C. Dong, H. Li, G. Mason, P. Montague:
\emph{The radical of a vertex operator algebra}. Definition der gradbewahrenden Mode $o(v)$.
\url{https://arxiv.org/abs/q-alg/9608022}.
\bibitem{barronhigherzhu} K. Barron, N. Vander Werf, J. Yang:
\emph{Higher level Zhu algebras and modules for vertex operator algebras}.
\url{https://arxiv.org/abs/1710.04767}.
\bibitem{entropyaction} A. H. Chamseddine, A. Connes, W. D. van Suijlekom:
\emph{Entropy and the spectral action}. Die Entropie der zweiten
Fermionquantisierung als Spektralaktion eines bereits gegebenen Diracoperators.
\url{https://arxiv.org/abs/1809.02944}.
\end{thebibliography}

\end{document}
